\subsection{Rank and Coordinate Models of Free Modules}

\begin{perspective} \label{per:dimension-as-field-case-of-rank} \hlt{Dimension as the Field Case of Rank}\\
An $\mathbb{F}$-vector space is a free $\mathbb{F}$-module. By Theorem \ref{thm:basis-length-well-defined}, finite bases have invariant cardinality; over general rings, rank invariance requires additional hypotheses.
\end{perspective}

\begin{theorem} \label{thm:free-module-rank-well-defined-commutative} \hlt{Rank Is Well Defined fo r Free Modules}\\
Let $A$ benonzero commutative ring with identity, $M$ be free $A$-module, $\operatorname{Bas}_A(M)$ denote set of $A$-bases of $M$.
\begin{enumerate}[label=(\roman*)]
\setlength{\itemsep}{0pt}
\item $B,B'\in\operatorname{Bas}_A(M)\Longrightarrow\abs{B}=\abs{B'}$.
\item $B\in\operatorname{Bas}_A(M)$, $S\subseteq M$, and $\langle\!\langle S\rangle\!\rangle_A=M\Longrightarrow\abs{B}\leq\abs{S}$.
\end{enumerate}
\begin{proof}
\begin{enumerate}[label=(\roman*)]
\setlength{\itemsep}{0pt}
\item Choose a maximal ideal $\mathcal{I}\lhd A$ by Theorem \ref{thm:maximal-ideals-exist}.\\
Put $K\defeq A/\mathcal{I}$, a field by Theorem \ref{thm:quotient-by-maximal-ideal-is-field}, $\mathcal{I}M\defeq\langle\!\langle im:i\in\mathcal{I},\ m\in M\rangle\!\rangle_A$, and $V\defeq M/\mathcal{I}M$.\\
Since $\mathcal{I}V=0$, $V$ is a $K$-vector space via $(a+\mathcal{I})(m+\mathcal{I}M)\defeq am+\mathcal{I}M$.
\item Let $B\in\operatorname{Bas}_A(M)$ and set $\overline{B}\defeq\{\overline{b}\defeq b+\mathcal{I}M:b\in B\}$.\\
In $B$-coordinates, $\mathcal{I}M=\{\sum\limits_{b\in F}i_bb:F\subseteq_{\operatorname{fin}}B,\ i_b\in\mathcal{I}\}$; hence $\overline{B}$ spans $V$.\\
If $\sum\limits_{j=1}^n(a_j+\mathcal{I})\overline{b_j}=0$ with distinct $b_j\in B$, then $\sum\limits_{j=1}^na_jb_j\in\mathcal{I}M$, so $a_j\in\mathcal{I}$ $\forall j$; hence $\overline{B}$ is $K$-independent.\\
If $\overline b=\overline{b'}$ with $b\neq b'$, then $b-b'\in\mathcal{I}M$, forcing the $B$-coefficient $1\in\mathcal{I}$, impossible. Thus $B\xrightarrow{\sim}\overline B$.
\item Therefore $\abs{B}=\abs{\overline B}=\dim_K V$ for every $B\in\operatorname{Bas}_A(M)$, so (i) follows from Theorem \ref{thm:arbitrary-bases-same-cardinality}.\\
If $\langle\!\langle S\rangle\!\rangle_A=M$, then $\overline S\defeq\{s+\mathcal{I}M:s\in S\}$ spans $V$, so $\abs{B}=\dim_K V\leq\abs{\overline S}\leq\abs S$.
\end{enumerate}
\hlt{Structural Commentary}: Modulo a maximal ideal, free module becomes a vector space with same basis size.\\
\hlt{Key Insight}: The field quotient $A/\mathcal{I}$ lets vector-space dimension certify module rank.
\end{proof}
\end{theorem}

\begin{definition} \label{def:rank-of-free-module} \hlt{Rank of a Free Module}\\
Let $A$ be a nonzero commutative ring with identity. The \hlt{rank} of a free
$A$-module $M$, denoted $\operatorname{rk}_A(M)$, is the cardinality of any
basis of $M$. The zero module has rank $0$.
\end{definition}

\begin{example} \label{ex:noncommutative-rank-failure-endomorphism-ring} \hlt{Why Commutativity Matters for Rank}\\
Let $V=\bigoplus_{k\geq0}\mathbb{F}e_k$ and $E=\LL(V)$. As a left $E$-module, $E=E\operatorname{id}_V$, so $\{\operatorname{id}_V\}$ is a basis.\\
For each $n\geq1$, define $\beta_s\in E$, $0\leq s<n$, by $\beta_s(e_{kn+t})=
\begin{cases}
e_k,&t=s,\\
0,&t\neq s,
\end{cases}$, where $k\geq0$, $0\leq t<n$.\\
If $\sum\limits_{s=0}^{n-1}\alpha_s\beta_s=0$, then evaluating at $e_{kn+t}$ gives $\alpha_t(e_k)=0$ for all $k,t$, hence $\alpha_t=0$ for all $t$.\\
Conversely, for $h\in E$, define $\alpha_s(e_k)\defeq h(e_{kn+s})$; then $h=\sum\limits_{s=0}^{n-1}\alpha_s\beta_s$.\\
Thus $\{\beta_0,\ldots,\beta_{n-1}\}$ is an $E$-basis of $E$, so $E$ has bases of all positive finite cardinalities.\\
\hlt{Structural Commentary}: Noncommutative scalar rings can make a free module have incompatible basis sizes.\\
\hlt{Key Insight}: The rank theorem is a commutative-ring theorem, not a formal consequence of freeness alone.
\end{example}

\begin{theorem} \label{thm:finite-support-functions-free-module} \hlt{Finite-Support Coordinate Modules}\\
Let $A$ be a commutative ring with identity and let $B$ be a set. Put $A^{(B)}\defeq\{f\in A^B:\abs{\operatorname{supp}(f)}<\infty\}$.\\
For $b\in B$, let $\delta_b\in A^{(B)}$ be given by $\delta_b(b)=1$ and $\delta_b(x)=0$ for $x\neq b$.\\
Then $A^{(B)}=\bigoplus\limits_{b\in B}A\delta_b$, so $\{\delta_b:b\in B\}$ is an $A$-basis. If $A\neq0$, then $\operatorname{rk}_A(A^{(B)})=\abs{B}$.
\begin{proof}
For $f,g\in A^{(B)}$ and $a\in A$, $\operatorname{supp}(f+g)\subseteq\operatorname{supp}(f)\cup\operatorname{supp}(g)$ and $\operatorname{supp}(af)\subseteq\operatorname{supp}(f)$, hence $A^{(B)}\leq A^B$.\\
For $f\in A^{(B)}$, $f=\sum\limits_{b\in\operatorname{supp}(f)}f(b)\delta_b$, so the $\delta_b$ span.
If $\sum\limits_{j=1}^{m}a_j\delta_{b_j}=0$ with $b_1,\ldots,b_m$ distinct, then evaluation at $b_i$ gives $a_i=0$ for every $i$. Thus the $\delta_b$ are independent, and the rank statement follows from the displayed basis.\\
\hlt{Structural Commentary}: Free modules are finite-support coordinate modules on their bases.\\
\hlt{Key Insight}: Finite support is the module-theoretic version of finite linear combination.
\end{proof}
\end{theorem}

\begin{proposition} \label{prop:free-module-coordinate-model-and-isomorphism} \hlt{Free Modules Are Coordinate Modules}\\
Let $A$ be a commutative ring with identity and let $M$ be free on a basis $B$.\\
Then the coordinate map $\Phi_B:M\to A^{(B)}$ is an isomorphism.\\
If $A\neq0$ and $M,N$ are free $A$-modules, then $M\cong N\Longleftrightarrow \operatorname{rk}_A(M)=\operatorname{rk}_A(N)$.
\begin{proof} For $m=\sum\limits_{b\in F}a_bb$, set $\Phi_B(m)=f_m$, with $f_m(b)=a_b$ for $b\in F$ and $f_m(b)=0$ for $b\notin F$.\\
Unique coordinates give $\Phi_B:M\xrightarrow{\sim}A^{(B)}$, with inverse $f\mapsto\sum\limits_{b\in\operatorname{supp}(f)}f(b)b$.
\begin{enumerate}[label=\roman*.]
\setlength{\itemsep}{0pt}
\item[$(\Rightarrow)$] If $\psi:M\xrightarrow{\sim}N$ and $C$ is a basis of $M$, then $\psi(C)$ is a basis of $N$, so $\operatorname{rk}_A(M)=\abs{C}=\abs{\psi(C)}=\operatorname{rk}_A(N)$.
\item[$(\Leftarrow)$] If $\operatorname{rk}_A(M)=\operatorname{rk}_A(N)$, choose bases $C,D$ with $\abs{C}=\abs{D}$.\\
Any bijection $C\to D$ extends, by Theorem \ref{thm:free-module-map-extension}, to $M\xrightarrow{\sim}N$.
\end{enumerate}
\noindent
\hlt{Structural Commentary}: A basis identifies a free module with a standard coordinate model.\\
\hlt{Key Insight}: Over a nonzero commutative ring, rank classifies free modules up to isomorphism.
\end{proof}
\end{proposition}

\begin{theorem} \label{thm:independent-sets-bounded-by-free-rank-domain} \hlt{Independent Sets Are Bounded by Rank over a Domain}\\
Let $A$ be integral domain with field of quotients $Q$, $M$ free over $A$, $S\subseteq M$ $A$-independent. Then $\abs{S}\leq\operatorname{rk}_A(M)$.
\begin{proof}
Choose an $A$-basis $B$ and identify $M=A^{(B)}\subseteq Q^{(B)}$.\\
If $\sum\limits_{i=1}^{n}q_is_i=0$ with $q_i\in Q$ and distinct $s_i\in S$, choose $d\in A\setminus\{0\}$ with $dq_i\in A$ for all $i$.\\
Then $\sum\limits_{i=1}^{n}(dq_i)s_i=0$ in $M$, so $dq_i=0$ by $A$-independence, hence $q_i=0$ for all $i$.\\
Thus $S$ is $Q$-independent in $Q^{(B)}$, and $\abs{S}\leq\dim_Q Q^{(B)}=\abs{B}=\operatorname{rk}_A(M)$.\\
\hlt{Structural Commentary}: The field of quotients lets ordinary dimension bound module independence.\\
\hlt{Key Insight}: Domains allow denominator clearing without creating false zero relations.
\end{proof}
\end{theorem}
